\documentclass[10pt,a4paper]{amsart}
\usepackage[margin=19mm]{geometry}
\usepackage{amsmath,amssymb,mathtools}
\usepackage[scr=rsfs,cal=euler]{mathalfa}
\usepackage{xfrac}
\usepackage[unicode,hidelinks,bookmarks=false]{hyperref}
\newcommand{\ie}{i.e.\ }
\newcommand{\cf}{cf.\ }
\newcommand{\resp}{respectively\ }
\newcommand{\Subsection}{\subsection}
\providecommand{\nobreakdash}{\nobreak\hbox{-}}
\setlength{\emergencystretch}{2em}
\multlinegap0pt
\usepackage{float}
\newtheorem{theorem}{Theorem}
\newtheorem{prop}[theorem]{Proposition}
\newtheorem{lemma}[theorem]{Lemma}
\newtheorem{definition}[theorem]{Definition}
\title[Lemma 4.5]{On Medial Quandle and Detecting Causality: the tangle lemma for medial quandles (Lemma 4.5)}
\author{Hongxu Chen}
\date{Source: arXiv:2411.04477v5 (CC BY 4.0)}
\begin{document}
\maketitle

\noindent\emph{Source and licence.} On Medial Quandle and Detecting Causality, by Hongxu Chen. Version arXiv:2411.04477v5
(CC BY 4.0), distributed by arXiv under the Creative Commons Attribution 4.0 International licence,
http://creativecommons.org/licenses/by/4.0/, as recorded in the arXiv licence metadata for this exact
version and shown on the abstract page https://arxiv.org/abs/2411.04477v5. Full licence notice:
\texttt{licenses/arxiv-2411.04477-CC-BY-4.0.txt}.\par\smallskip

\noindent\textit{Editorial scope.} Counted unit: the article's Lemma 4.5 (source0.tex lines 376--381, printed as Lemma 4.5): for a tangle $T_A$ given by a sub-diagram of $A_1$, any colouring of $T_A$ by a medial quandle forces $n=d$ and $m=e$. The article's complete original proof of it is delivered with it (source0.tex lines 383--425, one \texttt{proof} environment of 43 lines, adjacent to the statement, with no gap and no deferral). No mathematics is written, repaired or re-derived: the statement and the proof are the article's own lines, in the article's own notation, and no retained line is substituted or re-flowed. One bounded adaptation is declared and is the only difference from the source lines: the four image-inclusion commands inside the retained statement and proof (naming the four image files \texttt{TA.png}, \texttt{TAtransform.png}, \texttt{A1twostr.png} and \texttt{dissectofTA.png}) are omitted, so exactly four lines of the retained spans differ from the source and no other line does; the four figure environments themselves, with their placement, their captions and their labels, are carried verbatim, so every figure cross-reference of the retained lines still resolves. The package records these four omitted commands as its only adaptation, under a fidelity receipt bound to the counted proof, declares no asset dependency and ships no image. Retained uncounted as the source-local context the proof needs, each exactly the statement of the result it invokes and nothing else: the definition of the maps $f_{xy}$ (line 342, printed as Definition 4.2, source label \texttt{Dfxy}); the proposition listing the algebraic properties of those maps, with its five labelled items (lines 344--354, printed as Proposition 4.3, source label \texttt{propertyfxy}); and the theorem that the group they generate is abelian with automorphic elements (lines 364--368, printed as Theorem 4.4, source label \texttt{corofmapsareautomorphism}). Every cross-reference of every retained line resolves inside this document, and no retained line contains a citation, so the wrapper carries no bibliography and no bibliography entry.

Editorial formatting only, in addition: an AMS \texttt{amsart} preamble that restates the article's own declarations used by the retained lines -- the environments \texttt{definition}, \texttt{prop}, \texttt{theorem} and \texttt{lemma}, kept on one shared counter exactly as in the article but without the article's section-based prefix -- and the standard \texttt{float} package that the article itself loads, needed by the placement option of its misplaced-float figures. Because the section-based prefix is dropped, the retained environments print in this document as Definition 1, Proposition 2, Theorem 3 and Lemma 4 in order of appearance, the four retained figure environments print as Figure 1 to Figure 4, and the five retained displayed equations print as (1) to (5); the article prints the same items with the numbering of its own full document, which is not reproduced here because the items preceding them are not retained.

The complete native source of the article as distributed in the arXiv e-print for this exact version is bundled unchanged as \texttt{sources/168136/source0.tex}: it is byte identical to the e-print file \texttt{main.tex} except that the pipeline's mechanical concatenation prepends one header comment line naming it. The article was typeset with the standard \texttt{article} class; no publisher class or style file is shipped in the e-print and none is redistributed or used.
\begin{definition} \label{Dfxy} For any $x,y\in Q$, define $f_{xy}:Q\to Q$ by $f_{xy}(q)=(q*^{-1}x)*y$. Let $f_{x_2y_2}\circ f_{x_1y_1}$ denote the composition of two such functions. \end{definition} 

\begin{prop} \label{propertyfxy}\mbox{} 
Let $(Q,*)$ be a medial quandle. Then

\begin{enumerate}
    \item \label{a1} $[(a*x)*^{-1}y]*z=[(a*z)*^{-1}y]*x,~\forall a,x,y,z\in Q$.\\
\item \label{a2} $[(a*^{-1}x)*y]*^{-1}z=[(a*^{-1}z)*y]*^{-1}x,~\forall a,x,y,z\in Q$.\\ 
\item  \label{a3} $f_{xy}\circ f_{ab}=f_{ab}\circ f_{xy}$. 
    \item \label{a4} $f_{xy}(a*b)=f_{xy}(a)*f_{xy}(b)$ and $f_{xy}(a*^{-1}b)=f_{xy}(a)*^{-1}f_{xy}(b)$. 
    \item  \label{a5}$f_1\circ f_2\circ...\circ f_k=f_{1'}\circ f_{2'}\circ...\circ f_{k'}$, where $\{i\}\to\{i'\}$ is a permutation, and $f_i=f_{x_iy_i}$ for some $x_i,y_i\in Q$.

\end{enumerate} \end{prop}

\begin{theorem}
    
 \label{corofmapsareautomorphism}
The group generated by $\{f_{xy}~|~x,y\in Q\}$ is abelian and its elements are automorphisms on $Q$.
\end{theorem}

\begin{lemma} \label{lemma1}

Consider a tangle given by a sub-diagram of $A_1$ in Fig \ref{TA}, which is denoted by $T_A$.
       \begin{figure}[h]  \centering        \caption{The tangle $T_A$.} \label{TA}
          \end{figure}
\\For any coloring by medial quandle $(Q,*)$ on $T_A$, where $n,m,d,e,\in Q$ denote elements that color the two pairs of arcs of $T_A$ as illustrated in Fig \ref{TA}, we must have $n=d,m=e$. \end{lemma}

\begin{proof} \label{proflemma1}
     Note that from Figure \ref{TAtransf},  the following tangle is an equivalent form of $T_A$ by applying the third Reidemeister moves  : 
\begin{figure}[h]   \centering   
\caption{{Alternative diagram of $T_A$.} } \label{TAtransf}
\end{figure}
  
Let $(Q,*)$ be an arbitrary medial quandle. To analyze the coloring of $Q$ on $T_A$ We decompose $T_A$ into two tangles  We first analyze the coloring of $(Q,*)$ on two tangles shown in Figure \ref{fig:colortwostr}, called  $T_1$ and $T_2$, respectively. 
\begin{figure}[H]    \centering
  
 \caption{$T_1$(left) and $T_2$ (right). Recall the definition of $f_{xy}$ in Definition \ref{Dfxy}. }
    \label{fig:colortwostr} \end{figure}
We use letters denoting some elements in $Q$ to label arcs in Fig \ref{fig:colortwostr} to denote a quandle coloring on the structure. For both tangles, the elements that color the four arcs labeled by $x',y',n',m'$, can be expressed by $x,y,n,m$ by applying crossing relation : 
    

For structure \textcircled{1}:\begin{center}
    $n'=(n*^{-1}x)*y=f_{xy}(n)~~m'=(m*^{-1}x)*y=f_{xy}(m)$\\  $y'=y*x=(y*^{-1}y)*x=f_{yx}(y)$\\ $x'=x*^{-1}y'=x*^{-1}(y*x)=(x*^{-1}y)*(x*^{-1}x)=(x*^{-1}y)*x=f_{yx}(x)$ \end{center}

For structure \textcircled{2}: \begin{center}
     $a=f_{nm}(x)~~b=f_{nm}(y)$, so\\
$y'=b*a=f_{nm}(y)*f_{nm}(x)=f_{nm}(y*x)=f_{nm}((y*^{-1}y)*x)=f_{nm}\circ f_{yx}(y)$\\
$x'=a*^{-1}y'=f_{nm}(x)*^{-1}(f_{nm}\circ f_{yx}(y))=f_{nm}(x*^{-1}(y*x))=f_{nm}((x*^{-1}y)*x))=f_{nm}\circ f_{yx}(x)$ \end{center}

 Note that the diagram of $T_A$ in Fig \ref{TAtransf} can be viewed as the combination of $T_1$ and $T_2$:
  \begin{figure}[H] \centering  
  \caption{Decomposition of $T_A$.} \label{TAascombine} \end{figure}

  
  By applying the calculations from Figure \ref{fig:colortwostr} on the labels in Figure \ref{TAascombine}, we have \begin{equation}
       n'=f_{ji}(j),~j=f_{wv}\circ f_{ba}(b),~b=f_{yx}(y),~y=f_{ij}(v),~v=f_{xy}(n), \text{and} \label{lemma4.3 eq1}
  \end{equation}
  \begin{equation}
      m'=f_{ji}(i),~i=f_{wv}\circ f_{ba}(a),~a=f_{yx}(x),~x=f_{ij}(w),~w=f_{xy}(m). \label{lemma4.3 eq2}
  \end{equation}
   \\To simplify the expression for $j$ and $i$, we derived from equations \ref{lemma4.3 eq1} and \ref{lemma4.3 eq2} that 
  \begin{center}
       
   $b=f_{yx}\circ f_{ij}(v),~a=f_{yx}\circ f_{ij}(w)$.  \end{center} Let $g:=f_{yx}\circ f_{ij}$. By Theorem \ref{corofmapsareautomorphism}, $g$ is an automorphism on $Q$ and commutes with all $f_{xy}$-type maps. Therefore, 
   \begin{equation}
   j=f_{wv}\circ f_{ba}(b)=f_{wv}[(g(v)*^{-1}g(v))*g(w)]=f_{wv}(g(v*w))=g(f_{wv}(v*w))=g(v)=b. \end{equation} 
   \begin{equation} i=f_{wv}\circ f_{ba}(a)=f_{wv}[(g(w)*^{-1}g(v))*g(w)]=f_{wv}(g((w*^{-1}v)*w))=g(f_{wv}((w*^{-1}v)*w))=g(w)=a. \end{equation}
   Plugging in $j=b,~i=a$ to equations  \ref{lemma4.3 eq1} and \ref{lemma4.3 eq2} gives \begin{equation}
   n'=f_{ji}\circ f_{yx}\circ f_{ij}\circ f_{xy}(n), m'=f_{ji}\circ f_{yx}\circ f_{ij}\circ f_{xy}(m). \end{equation} Note that $f_{pq}\circ f_{qp}$ is the identity function for all $p,q\in Q(*)$ (recall Definition \ref{Dfxy}). By rearranging each individual functions which item \ref{a5} in Proposition \ref{propertyfxy} allows, we have\\
  $n'=f_{ji}\circ f_{ij}\circ f_{yx}\circ f_{xy}(n)=f_{ji}\circ f_{ij}(n)=n$ and $m'=f_{ji}\circ f_{ij}\circ f_{yx}\circ f_{xy}(m)=f_{ji}\circ f_{ij}(m)=m$, as desired. \end{proof}

\end{document}
